\section{Evaluation}
\label{sec:eval}

\subsection{Datasets}

% \begin{table}[h]
% \caption{Statistical information of the datasets.}
% \label{dataset}
% \centering
% \begin{tabular}{lcccc}
% \toprule
% \textbf{Statistics}  & \textbf{Agriculture} & \textbf{CS} & \textbf{Legal} & \textbf{Mix} \\
% \midrule
% Total Documents & 12  & 10  & 94  & 61  \\
% Total Tokens    & 2,017,886 & 2,306,535 & 5,081,069 & 619,009 \\
% \bottomrule
% \end{tabular}
% \end{table}

To conduct a comprehensive analysis of \model\, we select four datasets from the UltraDomain benchmark \cite{qian2024memorag}. The UltraDomain data is sourced from 428 college textbooks, covering 18 distinct domains, such as agriculture, social sciences and humanitie. From these, we chose the Agriculture, CS, Legal, and Mix datasets. The total number of tokens in each dataset ranges from 600,000 to 5,000,000, with detailed data shown in Table \ref{dataset}. A more detailed description is provided below.
\begin{itemize}
    \item \textbf{Agriculture}: Focuses on agricultural practice, covering topics like beekeeping, hive management, crop production, and disease prevention.
    
    \item  \textbf{CS}: Encompasses key areas of data science and software engineering, particularly machine learning and big data processing, featuring content on recommendation systems, classification algorithms, and real-time analytics using Spark.
    
    \item \textbf{Legal}: Centers around corporate legal practices, including materials on corporate restructuring, legal agreements, regulatory compliance, and governance, targeting the legal and financial sectors.
    
    \item \textbf{Mixed}: Offers a diverse range of literary, biographical, and philosophical texts, spanning cultural, historical, and philosophical studies.
\end{itemize}

To evaluate the effectiveness of RAG systems for more global sensemaking tasks, we merge all the text content from each dataset as context and follow the generation method used in \cite{edge2024local}. We task the LLM with identifying 5 potential users and 5 tasks for each user. For each (user, task) combination, the LLM is then required to generate 5 questions that necessitate an understanding of the entire corpus. In total, 125 questions are generated for each dataset.

\subsection{Baslines}
We compare \model\ with the following four baselines across all datasets:

\begin{enumerate}
    \item \textbf{Naive RAG} \cite{gao2023retrieval}: The most basic RAG system, which segments the text and stores the chunks in a vector database. Queries are vectorized and used to directly retrieve the most similar chunks from the database.
    
    \item \textbf{RQ-RAG} \cite{chan2024rq}: Utilizes the LLM to decompose the input query into several sub-queries. These sub-queries are designed to enhance search accuracy through explicit query rewriting, decomposition, and disambiguation.
    
    \item \textbf{HyDE} \cite{hyde}: Employs the LLM to generate a hypothetical document based on the input query. The hypothetical document is then used to retrieve relevant text chunks, which are subsequently used to formulate the final answer.
    
    \item \textbf{GraphRAG} \cite{edge2024local}: Leverages an LLM to extract entities and relationships as nodes and edges from the text and generates corresponding descriptions for them. It aggregates nodes into communities and generates a community report to capture global information. When faced with high-level queries, GraphRAG retrieves more comprehensive information by traversing the communities.
\end{enumerate}

\subsection{Settings}

For all experiments, we use the \href{https://github.com/gusye1234/nano-vectordb}{\textit{nano vector database}} to manage vector data. For the LLM used in the experiments, we utilize GPT-4o-mini. Chunk size across all datasets is set to 1200, and the gleaning parameter is fixed at 1 for both GraphRAG and \model\. As there is no established Ground Truth, we adopt the evaluation approach in \cite{edge2024local}. Each baseline is compared directly with \model\. In addition to the original four evaluation dimensions—Comprehensive, Diversity, Empowerment, and Directness—we introduce an additional dimension, Overall. The detailed criteria for each dimension are outlined as follows. The specific prompt is shown in Appendix x.

\begin{itemize}
    \item \textbf{Comprehensiveness}: How much detail does the answer provide to cover all aspects and details of the question?
    \item \textbf{Diversity}: How varied and rich is the answer in providing different perspectives and insights on the question?
    \item \textbf{Empowerment}: How well does the answer help the reader understand and make informed judgments about the topic?
    \item \textbf{Directness}: How specifically does the answer address the question?
    \item \textbf{Overall}: Considers the cumulative performance across all four dimensions (Comprehensiveness, Diversity, Empowerment, and Directness) to determine the best overall answer.
\end{itemize}

 The LLM directly compare two answers for each dimension and select the superior answer for each dimension. After determining the winning answer for each of the four dimensions, the LLM then combine the results across all dimensions to identify the overall better answer. To avoid any bias caused by the order of answers in the prompt during evaluation, we alternate the placement of each answer in the prompt and calculate win rates accordingly to obtain the final result.

\subsection{Results}
\begin{table}[h]
\centering
\caption{Win rates (\%) of Baselines vs \model\ across four datasets.}
\label{tab:performance}
\resizebox{\textwidth}{!}{
\begin{tabular}{@{}lcccccccc@{}}
\toprule
\textbf{}    & \multicolumn{2}{c}{\textbf{Agriculture}} & \multicolumn{2}{c}{\textbf{CS}} & \multicolumn{2}{c}{\textbf{Legal}} & \multicolumn{2}{c}{\textbf{Mix}} \\ 
\cmidrule(lr){2-3} \cmidrule(lr){4-5} \cmidrule(lr){6-7} \cmidrule(lr){8-9}
                      & NaiveRAG & \textbf{LightRAG} & NaiveRAG & \textbf{LightRAG} & NaiveRAG & \textbf{LightRAG} & NaiveRAG & \textbf{LightRAG} \\
\midrule
Comprehensiveness      & 32.69\%      & \underline{67.31\%}     & 35.44\%      & \underline{64.56\%}     & 19.05\%      & \underline{80.95\%}     & 36.36\%      & \underline{63.64\%}     \\
Diversity              & 24.09\%      & \underline{75.91\%}     & 35.24\%      & \underline{64.76\%}     & 10.98\%      & \underline{89.02\%}     & 30.76\%      & \underline{69.24\%}     \\
Empowerment            & 31.35\%      & \underline{68.65\%}     & 35.48\%      & \underline{64.52\%}     & 17.59\%      & \underline{82.41\%}     & 40.95\%      & \underline{59.05\%}     \\
Directness             & \underline{70.98\%}      & 29.02\%     & \underline{59.84\%}      & 40.16\%     & \underline{55.69\%}      & 44.31\%     & \underline{69.95\%}      & 30.05\%     \\
Overall                & 33.30\%      & \underline{66.70\%}     & 34.76\%      & \underline{65.24\%}     & 17.46\%      & \underline{82.54\%}     & 37.59\%      & \underline{62.40\%}     \\
\cmidrule(lr){2-3} \cmidrule(lr){4-5} \cmidrule(lr){6-7} \cmidrule(lr){8-9}
                      & RQ-RAG & \textbf{LightRAG} & RQ-RAG & \textbf{LightRAG} & RQ-RAG & \textbf{LightRAG} & RQ-RAG & \textbf{LightRAG} \\
\midrule
Comprehensiveness      & 32.05\%      & \underline{67.95\%}     & 39.30\%      & \underline{60.70\%}     & 18.57\%      & \underline{81.43\%}     & 38.89\%      & \underline{61.11\%}     \\
Diversity              & 29.44\%      & \underline{70.56\%}     & 38.71\%      & \underline{61.29\%}     & 15.14\%      & \underline{84.86\%}     & 28.50\%      & \underline{71.50\%}     \\
Empowerment            & 32.51\%      & \underline{67.49\%}     & 37.52\%      & \underline{62.48\%}     & 17.80\%      & \underline{82.20\%}     & 43.96\%      & \underline{56.04\%}     \\
Directness             & \underline{60.53\%}      & 39.47\%     & \underline{62.32\%}      & 37.68\%     & \underline{51.15\%}      & 48.85\%     & \underline{71.14\%}      & 28.86\%     \\
Overall         & 33.29\%      & \underline{66.71\%}     & 39.03\%      & \underline{60.97\%}     & 17.80\%      & \underline{82.20\%}     & 39.61\%      & \underline{60.39\%}     \\
\cmidrule(lr){2-3} \cmidrule(lr){4-5} \cmidrule(lr){6-7} \cmidrule(lr){8-9}
                      & HyDE & \textbf{LightRAG} & HyDE & \textbf{LightRAG} & HyDE & \textbf{LightRAG} & HyDE & \textbf{LightRAG} \\
\midrule
Comprehensiveness      & 24.39\%      & \underline{75.61\%}     & 36.49\%      & \underline{63.51\%}     & 27.68\%      & \underline{72.32\%}     & 42.17\%      & \underline{57.83\%}     \\
Diversity              & 24.96\%      & \underline{75.34\%}     & 37.41\%      & \underline{62.59\%}     & 18.79\%      & \underline{81.21\%}     & 30.88\%      & \underline{69.12\%}     \\
Empowerment            & 24.89\%      & \underline{75.11\%}     & 34.99\%      & \underline{65.01\%}     & 26.99\%      & \underline{73.01\%}     & \underline{45.61\%}      & 54.39\%     \\
Directness             & \underline{62.86\%}      & 37.14\%     & \underline{67.69\%}      & 32.31\%     & \underline{63.69\%}      & 36.31\%     & \underline{69.89\%}      & 30.11\%     \\
Overall         & 23.17\%      & \underline{76.83\%}     & 35.67\%      & \underline{64.33\%}     & 27.68\%      & \underline{72.32\%}     & 42.72\%      & \underline{57.28\%}     \\
\cmidrule(lr){2-3} \cmidrule(lr){4-5} \cmidrule(lr){6-7} \cmidrule(lr){8-9}
                      & GraphRAG & \textbf{LightRAG} & GraphRAG & \textbf{LightRAG} & GraphRAG & \textbf{LightRAG} & GraphRAG & \textbf{LightRAG} \\
\midrule
Comprehensiveness      & 45.56\%      & \underline{54.44\%}     & 45.98\%      & \underline{54.02\%}     &  47.13\%     & \underline{52.87\%}     & \underline{51.86\%}      & 48.14\%     \\
Diversity              & 19.65\%      & \underline{80.35\%}     & 39.64\%      & \underline{60.36\%}     &  25.55\%     & \underline{74.45\%}     & 35.87\%      & \underline{64.13\%}     \\
Empowerment            & 36.69\%      & \underline{63.31\%}     & 45.09\%      & \underline{54.91\%}     & 42.81\%      & \underline{57.19\%}     & \underline{52.94\%}      & 47.06\%     \\
Directness             & \underline{69.85\%}      & 30.15\%     & \underline{56.96\%}      & 43.04\%     & \underline{58.11\%}      & 41.89\%     & \underline{55.44\%}      & 44.56\%     \\
Overall                & 43.62\%      & \underline{56.38\%}     & 45.98\%      & \underline{54.02\%}     & 45.70\%      & \underline{54.30\%}     & \underline{51.86\%}      & 48.14\%     \\
\bottomrule
\end{tabular}
}
\end{table}

To comprehensively evaluate the performance of \model\, we compare it against baselines across four datasets, with results reported in Table \ref{tab:performance}. These results are measured across five dimensions. Due to inherent trade-offs between some of these dimensions, it is unlikely that any system would outperform in all aspects simultaneously \cite{edge2024local}. By examining the data, three major conclusions can be drawn:

\textbf{Superiority of Graph-based Retrieval in Large-Scale Token Queries}: 
When dealing with large token counts and complex queries requiring comprehensive consideration of the dataset's background, Graph-based RAG systems like \model\ consistently outperform chunk-based retrieval methods (NaiveRAG, HyDE, and RQRAG). This performance gap becomes particularly evident as the dataset size increases. For instance, in the smallest dataset (Mix), the performance difference between chunk-based methods and \model\ is relatively modest. However, in the largest dataset (Legal), the gap widens significantly, with chunk-based methods only achieving around 20\% \% win rates overall compared to \model\'s dominance. This trend highlights the scalability advantage of Graph-based retrieval, especially in handling more extensive and intricate datasets.
    
\textbf{Consistent Leadership in the Diversity Dimension}: 
Across comparisons with chunk-based baselines (NaiveRAG, HyDE, RQRAG), \model\ demonstrates superior performance in three key dimensions: Comprehensiveness, Diversity, and Empowerment, with its advantage in the Diversity dimension being particularly pronounced. Even when compared with the Graph-based GraphRAG system, \model\ shows a noticeable edge in Diversity, especially in the larger Legal dataset. Although \model\ slightly underperforms in the Overall dimension on the Mix dataset compared to GraphRAG, its consistent lead in Diversity suggests that it is more effective at generating a broader range of responses, particularly in scenarios where diverse content is critical.
    
\textbf{\model's Superiority over GraphRAG}: 
Although both \model\ and GraphRAG leverage graph-based retrieval mechanisms, \model\ consistently outperforms GraphRAG across multiple dimensions, particularly in larger and more complex datasets. In the Agriculture, CS, and Legal datasets, which each contain millions of tokens, \model\ achieved a clear and consistent lead, comprehensively surpassing GraphRAG. In the smallest dataset, Mix, while \model\ showed a slight disadvantage in the Overall dimension, its performance remained within a comparable range. More importantly, \model\ still outperformed GraphRAG in the Diversity dimension, further demonstrating \model\'s strong global control over the datasets, reinforcing its advantage in diverse and complex data environments. By generating keywords during both graph creation and query processing, \model\ ensures that the graph remains scalable, efficiently handling a larger scope of data while maintaining performance.

\subsection{Compare Three modes of LightRAG}

\begin{table}[h]
\centering
\caption{}
\label{tab:ablation}
\resizebox{\textwidth}{!}{
\begin{tabular}{@{}lcccccccc@{}}
\toprule
\textbf{}    & \multicolumn{2}{c}{\textbf{Agriculture}} & \multicolumn{2}{c}{\textbf{CS}} & \multicolumn{2}{c}{\textbf{Legal}} & \multicolumn{2}{c}{\textbf{Mix}} \\ 
\cmidrule(lr){2-3} \cmidrule(lr){4-5} \cmidrule(lr){6-7} \cmidrule(lr){8-9}
                      & NaiveRAG & \textbf{Hybird} & NaiveRAG & \textbf{Hybird} & NaiveRAG & \textbf{Hybird} & NaiveRAG & \textbf{Hybird} \\
\midrule
Comprehensiveness      & 32.69\%      & \underline{67.31\%}     & 35.44\%      & \underline{64.56\%}     & 19.05\%      & \underline{80.95\%}     & 36.36\%      & \underline{63.64\%}     \\
Diversity              & 24.09\%      & \underline{75.91\%}     & 35.24\%      & \underline{64.76\%}     & 10.98\%      & \underline{89.02\%}     & 30.76\%      & \underline{69.24\%}     \\
Empowerment            & 31.35\%      & \underline{68.65\%}     & 35.48\%      & \underline{64.52\%}     & 17.59\%      & \underline{82.41\%}     & 40.95\%      & \underline{59.05\%}     \\
Directness             & \underline{70.98\%}      & 29.02\%     & \underline{59.84\%}      & 40.16\%     & \underline{55.69\%}      & 44.31\%     & \underline{69.95\%}      & 30.05\%     \\
Overall                & 33.30\%      & \underline{66.70\%}     & 34.76\%      & \underline{65.24\%}     & 17.46\%      & \underline{82.54\%}     & 37.59\%      & \underline{62.40\%}     \\
\cmidrule(lr){2-3} \cmidrule(lr){4-5} \cmidrule(lr){6-7} \cmidrule(lr){8-9}
                      & NaiveRAG & \textbf{Local} & NaiveRAG & \textbf{Local} & NaiveRAG & \textbf{Local} & NaiveRAG & \textbf{Local} \\
\midrule
Comprehensiveness      & 35.79\%      & \underline{64.21\%}     & 42.98\%      & \underline{57.02\%}     & 22.60\%      & \underline{77.40\%}     & 40.69\%      & \underline{59.31\%}     \\
Diversity              & 26.86\%      & \underline{73.14\%}     & 35.09\%      & \underline{64.91\%}     & 16.09\%      & \underline{83.91\%}     & 37.15\%      & \underline{62.84\%}     \\
Empowerment            & 35.02\%      & \underline{64.98\%}     & 42.98\%      & \underline{57.02\%}     & 23.58\%      & \underline{76.42\%}     & 48.26\%      & \underline{51.74\%}     \\
Directness             & \underline{67.90\%}      & 32.10\%     & \underline{70.18\%}      & 29.82\%     & \underline{51.57\%}      & 48.43\%     & \underline{63.96\%}      & 36.04\%     \\
Overall                & 35.33\%      & \underline{64.67\%}     & 42.98\%      & \underline{57.02\%}     & 21.92\%      & \underline{78.08\%}     & 41.39\%      & \underline{58.61\%}     \\
\cmidrule(lr){2-3} \cmidrule(lr){4-5} \cmidrule(lr){6-7} \cmidrule(lr){8-9}
                      & NaiveRAG & \textbf{Global} & NaiveRAG & \textbf{Global} & NaiveRAG & \textbf{Global} & NaiveRAG & \textbf{Global} \\
\midrule
Comprehensiveness      & 36.31\%      & \underline{63.69\%}     & 41.82\%      & \underline{58.18\%}     & 18.92\%      & \underline{81.08\%}     & 35.92\%      & \underline{64.08\%}     \\
Diversity              & 29.07\%      & \underline{70.93\%}     & 40.29\%      & \underline{59.71\%}     & 13.06\%      & \underline{86.94\%}     & 32.34\%      & \underline{67.66\%}     \\
Empowerment            & 34.87\%      & \underline{65.13\%}     & 41.51\%      & \underline{58.49\%}     & 17.73\%      & \underline{82.27\%}     & 35.98\%      & \underline{64.02\%}     \\
Directness             & \underline{71.19\%}      & 28.81\%     & \underline{60.17\%}      & 39.83\%     & \underline{57.17\%}      & 42.83\%     & \underline{69.41\%}      & 30.59\%     \\
Overall                & 35.02\%      & \underline{64.98\%}     & 42.00\%      & \underline{58.00\%}     & 19.00\%      & \underline{81.00\%}     & 35.92\%      & \underline{64.08\%}     \\
\bottomrule
\end{tabular}
}
\end{table}

To explore the performance of different \model\ modes across various dimensions, we compared three modes of \model\ with NaiveRAG across four datasets, and the results are presented in Table \ref{tab:ablation}. 

\textbf{Local Mode}: From an overall perspective, Local mode is at a disadvantage. This is mainly due to its retrieval mechanism, which focuses on entities and their directly related relationships. While this allows for deeper exploration of specific content, it struggles to address queries that require more extensive or broader information. As a result, Local mode often falls short in empowerment when compared to other modes. Because the in-depth information focus on specific content does not always effectively answer queries that require broader and more comprehensive insights. However, this focus retrieval method enables it to provide more detailed and in-depth content, which explains its strong performance in the directness dimension.

\textbf{Global Mode}: Global mode leverages a keyword-based retrieval method that allows it to capture a broader range of content by focusing on relationships between entities, rather than drilling down into specific entities. This approach provides a clear advantage in the comprehensiveness dimension, as it is able to gather more extensive and varied information. However, the trade-off is a lack of depth in examining specific events or entities, which can limit its ability to provide highly detailed insights. As a result, Global mode may struggle with tasks that demand more precise, detailed answers.

\textbf{Hybird Mode}: Hybird mode combines the strengths of both Local and Global modes, retrieving a broader set of relationships while simultaneously conducting an in-depth exploration of specific entities. This dual approach ensures both breadth in the retrieval process and depth in the analysis, offering a comprehensive view of the data. As a result, Hybird mode achieves a balanced performance across multiple dimensions.

\subsection{Impact of Original Text}
\subsection{Impact of Topk}
\begin{table}[h]
\centering
\caption{}
\resizebox{\textwidth}{!}{
\begin{tabular}{@{}lcccccccc@{}}
\toprule
\textbf{}    & \multicolumn{2}{c}{\textbf{Agriculture}} & \multicolumn{2}{c}{\textbf{CS}} & \multicolumn{2}{c}{\textbf{Legal}} & \multicolumn{2}{c}{\textbf{Mix}} \\ 
\cmidrule(lr){2-3} \cmidrule(lr){4-5} \cmidrule(lr){6-7} \cmidrule(lr){8-9}
                      & NaiveRAG & \textbf{Hybird} & NaiveRAG & \textbf{Hybird} & NaiveRAG & \textbf{Hybird} & NaiveRAG & \textbf{Hybird} \\
\midrule
Comprehensiveness      & 32.69\%      & \underline{67.31\%}     & 35.44\%      & \underline{64.56\%}     & 19.05\%      & \underline{80.95\%}     & 36.36\%      & \underline{63.64\%}     \\
Diversity              & 24.09\%      & \underline{75.91\%}     & 35.24\%      & \underline{64.76\%}     & 10.98\%      & \underline{89.02\%}     & 30.76\%      & \underline{69.24\%}     \\
Empowerment            & 31.35\%      & \underline{68.65\%}     & 35.48\%      & \underline{64.52\%}     & 17.59\%      & \underline{82.41\%}     & 40.95\%      & \underline{59.05\%}     \\
Directness             & \underline{70.98\%}      & 29.02\%     & \underline{59.84\%}      & 40.16\%     & \underline{55.69\%}      & 44.31\%     & \underline{69.95\%}      & 30.05\%     \\
Overall                & 33.30\%      & \underline{66.70\%}     & 34.76\%      & \underline{65.24\%}     & 17.46\%      & \underline{82.54\%}     & 37.59\%      & \underline{62.40\%}     \\
\cmidrule(lr){2-3} \cmidrule(lr){4-5} \cmidrule(lr){6-7} \cmidrule(lr){8-9}
                      & NaiveRAG & \textbf{Hybird-wo} & NaiveRAG & \textbf{Hybird-wo} & NaiveRAG & \textbf{Hybird-wo} & NaiveRAG & \textbf{Hybird-wo} \\
\midrule
Comprehensiveness      & 25.23\%      & \underline{74.77\%}     & 40.20\%      & \underline{59.80\%}     & 16.40\%      & \underline{83.60\%}     & 44.00\%      & \underline{56.00\%}     \\
Diversity              & 26.94\%      & \underline{73.06\%}     & 45.36\%      & \underline{54.64\%}     & 13.15\%      & \underline{86.84\%}     & 26.00\%      & \underline{74.00\%}     \\
Empowerment            & 31.01\%      & \underline{68.99\%}     & 42.61\%      & \underline{57.39\%}     & 18.55\%      & \underline{81.45\%}     & 45.33\%      & \underline{54.67\%}     \\
Directness             & \underline{61.94\%}      & 38.06\%     & \underline{62.07\%}      & 37.93\%     & \underline{54.51\%}      & 45.49\%     & \underline{63.33\%}      & 36.67\%     \\
Overall                & 25.51\%      & \underline{74.49\%}     & 40.20\%      & \underline{59.80\%}     & 17.12\%      & \underline{82.88\%}     & 44.00\%      & \underline{56.00\%}     \\
\bottomrule
\end{tabular}
}
\end{table}

\begin{figure}[h]
\centering
\includegraphics[width=1\textwidth]{figs/evaluation_prompt.png}
\caption{}
\label{fig:evaluation_prompt}
\end{figure}

\subsection{Case Study}

\begin{table}[h]
\centering
\caption{}
\label{tab:metrics}
\resizebox{\textwidth}{!}{
\begin{tabular}{|l|p{15cm}|}
\hline
\textbf{Query} & Which methods can normalize the feature values for effective machine learning? \\ \hline
\textbf{GraphRAG} & \textbf{Key Metrics for Evaluating Movie Recommendation Systems}:  
Evaluating the effectiveness of movie recommendation systems is crucial for ensuring user satisfaction and improving the overall performance of these systems. Several metrics are widely recognized for their ability to provide insights into how well a recommendation system performs. Below are the most informative metrics:  
1. \textbf{Precision} [...]  
2. \textbf{Recall} [...]  
3. \textbf{F1 Score} [...]  
4. \textbf{Mean Average Precision (MAP)} [...]  
5. \textbf{Root Mean Squared Error (RMSE)} [...]  
6. \textbf{User Satisfaction Metrics} [...]  
\textbf{Conclusion}: In summary, metrics such as precision, recall, F1 score, MAP, RMSE, and user satisfaction are critical for evaluating the performance of movie recommendation systems. Each metric provides unique insights that can help developers and researchers refine their algorithms and enhance user experiences. By focusing on these metrics, stakeholders can ensure that their recommendation systems effectively meet user preferences and expectations. \\ \hline
\textbf{LightRAG} & \textbf{Key Metrics for Evaluating Movie Recommendation Systems}:  
Evaluating movie recommendation systems involves multiple metrics that assess the effectiveness and accuracy of the recommendations provided to users. Below are some of the most informative metrics commonly used in this context, derived from the data and concepts surrounding recommendation systems.  
1. \textbf{Mean Average Precision at K (MAPK)} [...]  
2. \textbf{Precision and Recall} [...]  
3. \textbf{Root Mean Squared Error (RMSE) and Mean Squared Error (MSE)} [...]  
4. \textbf{Area Under the Curve (AUC)} [...]  
5. \textbf{F-Measure} [...]  
6. \textbf{User Engagement Metrics} [...]  
\textbf{Conclusion}: A comprehensive evaluation of movie recommendation systems employs a range of metrics that capture both the accuracy of the recommendations and user satisfaction. By effectively leveraging these metrics, system designers can optimize recommendation engines to deliver more personalized and relevant experiences to users. \\ \hline
\textbf{LLM Decision} & \textbf{Comprehensiveness}:  
"Winner": "Answer 2 (LightRAG)",  \newline

"Explanation": "Answer 2 provides a broader array of metrics including MAPK, AUC, and user engagement metrics alongside precision, recall, RMSE, and F-Measure. This gives a fuller view of how recommendation systems can be evaluated. Answer 1, while detailed, is more focused on traditional metrics and lacks the inclusion of engagement metrics which are crucial for understanding user interaction." \newline

\textbf{Diversity}:  
"Winner": "Answer 2 (LightRAG)", \newline

"Explanation": "Answer 2 not only covers a wide variety of metrics but also includes nuanced explanations of how some metrics interrelate and differ from one another, like the inclusion of both RMSE and MSE, as well as the details behind AUC. In contrast, Answer 1 sticks primarily to standard metrics without much exploration of potential nuances." \newline

\textbf{Empowerment}:  
"Winner": "Answer 2 (LightRAG)",  \newline

"Explanation": "Answer 2 empowers the reader more effectively by detailing how each metric functions and its importance in evaluating recommendation systems. By providing context such as the trade-offs between precision and recall and emphasizing user engagement metrics, it enables readers to make more informed judgments and understand the implications of different metrics. Answer 1 is more straightforward but lacks the depth of insight regarding why these metrics matter." \newline

% \textbf{Directness}:  
% "Winner": "Answer 1 (GraphRAG)",  \newline

% "Explanation": "Answer 1 presents each metric systematically with clear definitions and implications. It directly addresses the question with straightforward definitions and practical implications for each listed metric. While Answer 2 is also direct, its additional elaborations sometimes lead to a slightly more complex presentation, making it less immediately accessible." \newline

\textbf{Overall Winner}:  
"Winner": "Answer 2 (LightRAG)",  \newline

"Explanation": "While Answer 1 is more direct and systematic, Answer 2 excels in comprehensiveness, diversity, and empowerment. It provides a richer exploration of the topic, including insights into user engagement and nuanced differences between metrics. This depth and breadth make it more informative for readers seeking to thoroughly understand the evaluation of movie recommendation systems." \\ \hline

\end{tabular}
}
\end{table}


\subsection{Complexity Analysis}

\begin{table}[h]
\centering
\begin{tabular}{|l|l|l|l|}
\hline
\multicolumn{2}{|c|}{\textbf{阶段}}          & \multicolumn{2}{c|}{\textbf{方法}}                     \\ \hline
\textbf{}          & \textbf{} & \textbf{GraphRAG}                                      & \textbf{LightRAG}                              \\ \hline
\textbf{检索阶段}   & Token开销  & 610,000 tokens（610 communities × 1,000 tokens/community） & < 100 tokens                                   \\ \hline
                    & API调用次数 & 数百次（遍历610个community）                              & 1次调用                                       \\ \hline
\textbf{新增文本阶段} & Token开销  & 1,399 * 2 * 5,000 tokens（重新生成社区报告）                & 无显著额外token开销                           \\ \hline
                    & API调用次数 & 数百次（遍历重新生成的community）                          & 1次调用                                       \\ \hline
\end{tabular}
\caption{GraphRAG与LightRAG在检索阶段和新增文本阶段的开销对比}
\end{table}


On the legal dataset during the retrieval phase, GraphRAG generates and utilizes a total of 1,399 communities, of which 610 level-2 communities are actively used for retrieval. Each community report contains an average of 1,000 tokens, resulting in a total token consumption of 610,000 tokens (610 communities × 1,000 tokens per community). Additionally, due to the need to traverse each community individually, hundreds of API calls are required, significantly increasing the retrieval overhead. In contrast, LightRAG optimizes this process by using less than 100 tokens for keyword generation and retrieval. And only requires a single API call for the entire retrieval process, further reducing the computational and time overhead compared to GraphRAG.


For newly added text, the overhead for entity and relationship extraction is similar for both GraphRAG and LightRAG. However, due to GraphRAG's need to reconstruct communities and generate community reports, there is a significant difference in the subsequent processing stages. Taking the legal dataset as an example, for a newly added dataset of the same size, GraphRAG first needs to break down the existing community structure, incorporate the new entities and relationships, and then regenerate the community structure. The token consumption for generating each community report is approximately 5,000 tokens. In the original legal dataset, there are 1,399 communities, so for the new dataset, GraphRAG would need to reconstruct both the original and new community reports, leading to a total consumption of about 1,399 * 2 * 5,000 tokens. In contrast, LightRAG only needs to integrate the newly extracted entities and relationships into the original graph without reconstructing the entire community structure. This makes LightRAG more efficient in handling incremental updates.

